\section*{Justificación}
\addcontentsline{toc}{section}{Justificación}

Automatizar el diseño de circuitos analógicos tiene valor práctico porque su ciclo de propuesta, simulación y ajuste descansa en el conocimiento experto y no puede seguirse como un procedimiento mecánico. Los modelos de lenguaje que generan diseños a partir de lenguaje natural simplifican la entrada al proceso, pero producen diseños eléctricamente inconsistentes que requieren corrección manual o, cuando los sistemas verifican internamente, lo hacen a nivel de transistor o mediante reglas eléctricas, no mediante simulación SPICE dentro del lazo de diseño.

En la práctica, a partir de una descripción en lenguaje natural el sistema genera un netlist con componentes comerciales, lo simula en ngspice \cite{ngspice_manual}, extrae sus métricas eléctricas y ajusta los valores de los componentes de forma iterativa hasta que el diseño cumpla las metas especificadas. El ciclo de ajuste lo ejecuta el propio ecosistema; el usuario recibe el netlist simulado junto con las métricas y el historial de iteraciones.

Los destinatarios directos son ingenieros y estudiantes que trabajen con circuitos analógicos CA/CD construidos a partir de componentes comerciales con modelo SPICE disponible ---resistores, capacitores, inductores, fuentes, diodos, transistores y amplificadores operacionales de catálogo--- y que requieran una herramienta que no solo genere sino que también valide mediante simulación, sin depender de un nodo tecnológico propietario ni de herramientas de licencia costosa.

Lo que se propone es trabajar con componentes comerciales y, a la vez, usar la simulación SPICE como señal de ajuste dentro del lazo de diseño, coordinado por varios agentes con un curador. Se apoya en herramientas de uso libre ---LangGraph \cite{langgraph}, PySpice \cite{pyspice} y ngspice \cite{ngspice_manual}---, desplegables en la nube y sin dependencia de licencias. Desde la ingeniería en computación.

Este proyecto implementa un ecosistema multiagente, que integra tareas mediante procesamiento paralelo y obedece el comportamiento eléctrico, siendo una herramienta mas acertada en las generación y validación de circuitos mediante el uso de LLM’s.

